\documentclass{ibetest}
\renewcommand{\arraystretch}{1.3}
\begin{document}
\testheader
\vspace{14mm}
\begin{center}
  {\sffamily\bfseries\Huge Thermo-S1 \,$\cdot$\, Part 1}\\[6pt]
  {\sffamily\bfseries\LARGE Answer keys and marking schemes}\\[8pt]
  {\large Progress Tests 1.A $\to$ 3.C}\\[3pt]
  {\normalsize Semester 1 \,$\cdot$\, Academic Year 2026--2027}
\end{center}
\vspace{9mm}

{\sffamily\bfseries How to mark yourself}
\begin{itemize}[leftmargin=7mm, itemsep=3pt]
  \item Model answers are in \textcolor{keyblue}{blue}, inside the same frames as on the
        paper. Your wording may differ; what matters is the content listed in the
        \emph{marking scheme} below each exercise.
  \item Each circled item is worth the points shown. A numerical result without its unit,
        or with a wrong unit symbol, does not earn its item.
  \item Tick-box questions: 1 mark each. The ``Why'' line explains the answer and the trap
        behind the most tempting wrong option.
  \item Red boxes (\emph{Examiner's tip}) flag common traps, including the ones that cost
        marks in Test~1.
  \item Several boxes of the handout are blank (for example Boxes 4--6 and 8--18). The keys
        give the standard definitions; accept the wording used in class.
\end{itemize}

\vspace{7mm}
{\sffamily\bfseries\large Lessons from Test~1, applied to thermodynamics}\par\vspace{3mm}
\begin{center}
\small
\begin{tabular}{@{}>{\raggedright\arraybackslash}p{58mm} >{\raggedright\arraybackslash}p{92mm}@{}}
\hline
\textbf{Test~1 (electricity)} & \textbf{In thermodynamics}\\
\hline
``KV'': the marker wrote ``big K stands for kelvins''. &
  Temperatures are in kelvins: ``300 K'', never ``300\,$^\circ$K''. The kilo prefix stays
  lowercase: kPa.\\ \hline
The unit of current is the ampere (A), an SI base unit, not ``C\,s$^{-1}$''. &
  The kelvin (K) and the mole (mol) are SI base units. The pascal is derived:
  $\mathrm{Pa} = \mathrm{kg\,m^{-1}\,s^{-2}}$, and $R$ is in $\mathrm{J\,K^{-1}\,mol^{-1}}$.\\
  \hline
Read the circuit diagram exactly as drawn (nodes, branches). &
  Read state diagrams exactly: the fixed variable (isotherm, isobar, isochore) decides which
  partial derivative a slope is.\\ \hline
Isolate the unknown first, carry units, handle powers of ten separately
($v = i/neS$, then $3.1\times10^{-5}$~m/s). &
  Write $p = nRT/V$ first, then substitute SI values: $1\ \mathrm L = 10^{-3}\ \mathrm{m^3}$,
  $1\ \mathrm{bar} = 10^{5}\ \mathrm{Pa}$, $1\ \mathrm{cm^2} = 10^{-4}\ \mathrm{m^2}$.\\ \hline
\end{tabular}
\end{center}

\vspace{6mm}
\begin{tcolorbox}[enhanced, colback=black!3, colframe=black!45, boxrule=0.4pt, arc=1mm,
  fontupper=\small]
\textbf{Numbers without a calculator.} The papers give $R = 8.3\ \mathrm{J\,K^{-1}\,mol^{-1}}$
and $RT\approx2.5\times10^{3}\ \mathrm{J\,mol^{-1}}$ at 300~K. With the exact
$R = 8.314\ \mathrm{J\,K^{-1}\,mol^{-1}}$, the answers change by less than 1\,\%.
\end{tcolorbox}
\end{document}
